%% intro.tex
%%
%% Copyright (C) 2000-2018 Simone Piccardi.  Permission is granted to
%% copy, distribute and/or modify this document under the terms of the GNU Free
%% Documentation License, Version 1.1 or any later version published by the
%% Free Software Foundation; with the Invariant Sections being "Un preambolo",
%% with no Front-Cover Texts, and with no Back-Cover Texts.  A copy of the
%% license is included in the section entitled "GNU Free Documentation
%% License".
%%

\chapter{L'architettura del sistema}
\label{cha:intro_unix}

In questo primo capitolo sarà fatta un'introduzione ai concetti generali su
cui è basato un sistema operativo di tipo Unix come GNU/Linux, in questo modo
potremo fornire una base di comprensione mirata a sottolineare le peculiarità
del sistema che sono più rilevanti per quello che riguarda la programmazione.

Dopo un'introduzione sulle caratteristiche principali di un sistema di tipo
Unix passeremo ad illustrare alcuni dei concetti base dell'architettura di
GNU/Linux (che sono comunque comuni a tutti i sistemi \textit{unix-like}) ed
introdurremo alcuni degli standard principali a cui viene fatto riferimento.


\section{Una panoramica}
\label{sec:intro_unix_struct}

In questa prima sezione faremo una breve panoramica sull'architettura di un
sistema operativo di tipo Unix, come GNU/Linux, e della relazione fra le varie
parti che lo compongono. Chi avesse già una conoscenza di questa materia può
tranquillamente saltare questa sezione.

\subsection{Concetti base}
\label{sec:intro_base_concept}

Il concetto principale su cui è basata l'architettura di un sistema unix-like
è quello di un nucleo del sistema, il cosiddetto \textit{kernel} (nel nostro
caso Linux) a cui si demanda la gestione delle risorse della propria macchina
(la CPU, la memoria, le periferiche) mentre tutto il resto, quindi anche la
parte che prevede l'interazione con l'utente, dev'essere realizzato tramite
programmi eseguiti dal kernel, che accedano alle risorse tramite opportune
richieste a quest'ultimo.

Fin dai suoi albori Unix nasce come sistema operativo \textit{multitasking},
cioè in grado di eseguire contemporaneamente più programmi, e multiutente, in
cui è possibile che più utenti siano connessi ad una macchina eseguendo più
programmi ``\textsl{in contemporanea}''. In realtà, almeno per le macchine a
processore singolo, i programmi vengono semplicemente eseguiti uno alla volta
in una opportuna \textsl{rotazione}.\footnote{anche se oggi, con la presenza
  di sistemi multiprocessore, si possono avere più processi eseguiti in
  contemporanea, il concetto di ``\textsl{rotazione}'' resta comunque valido,
  dato che in genere il numero di processi da eseguire eccede il numero dei
  precessori disponibili. }

% Questa e` una distinzione essenziale da capire,
%specie nei confronti dei sistemi operativi successivi, nati per i personal
%computer (e quindi per un uso personale), sui quali l'hardware (allora
%limitato) non consentiva la realizzazione di un sistema evoluto come uno unix.

I kernel Unix più recenti, come Linux, sono realizzati sfruttando alcune
caratteristiche dei processori moderni come la gestione hardware della memoria
e la modalità protetta. In sostanza con i processori moderni si può
disabilitare temporaneamente l'uso di certe istruzioni e l'accesso a certe
zone di memoria fisica.  Quello che succede è che il kernel è il solo
programma ad essere eseguito in modalità privilegiata, con il completo accesso
a tutte le risorse della macchina, mentre i programmi normali vengono eseguiti
in modalità protetta senza accesso diretto alle risorse.  Uno schema
elementare della struttura del sistema è riportato in
fig.~\ref{fig:intro_sys_struct}.

\begin{figure}[htb]
  \centering
  \includegraphics[width=11cm]{img/struct_sys}
  % \begin{tikzpicture}
  %   \filldraw[fill=black!20] (0,0) rectangle (7.5,1);
  %   \draw (3.75,0.5) node {\textsl{System Call Interface}};
  %   \filldraw[fill=black!35] (0,1) rectangle (7.5,4);
  %   \draw (3.75,2.5) node {\huge{\textsf{kernel}}};
  %   \filldraw[fill=black!20] (0,4) rectangle (2.5,5);
  %   \draw (1.25,4.5) node {\textsf{scheduler}};
  %   \filldraw[fill=black!20] (2.5,4) rectangle (5,5);
  %   \draw (3.75,4.5) node {\textsf{VM}};
  %   \filldraw[fill=black!20] (5,4) rectangle (7.5,5);
  %   \draw (6.25,4.5) node {\textsf{driver}};

  %   \draw (1.25,7) node(cpu) [ellipse,draw] {\textsf{CPU}};
  %   \draw (3.75,7) node(mem) [ellipse,draw] {\textsf{memoria}};
  %   \draw (6.25,7) node(disk) [ellipse,draw] {\textsf{disco}};

  %   \draw[<->] (cpu) -- (1.25,5);
  %   \draw[<->] (mem) -- (3.75,5);
  %   \draw[<->] (disk) -- (6.25,5);

  %   \draw (7.5,0) node [anchor=base west] {\textit{kernel space}};
  %   \draw (7.5,-1) node [anchor=west] {\textit{user space}};

  %   \draw (-1,-0.5) -- (8.5, -0.5);

  %   \draw (0,-2) rectangle (7.5,-1);
  %   \draw (3.75, -1.5) node {\textsl{GNU C Library}};
  %   \draw[->] (1.25,-1) -- (1.25,0);
  %   \draw[->] (3.75,-1) -- (3.75,0);
  %   \draw[->] (6.25,-1) -- (6.25,0);

  %   \draw (1.25,-3) node(proc1) [rectangle,draw] {\textsf{processo}};
  %   \draw (3.75,-3) node(proc2) [rectangle,draw] {\textsf{processo}};
  %   \draw (6.25,-3) node(proc3) [rectangle,draw] {\textsf{processo}};

  %   \draw[->] (1.25,-2) -- (proc1);
  %   \draw[->] (3.75,-2) -- (proc2);
  %   \draw[->] (6.25,-2) -- (proc3);
  % \end{tikzpicture}
  \caption{Schema di massima della struttura di interazione fra processi,
    kernel e dispositivi in Linux.}
  \label{fig:intro_sys_struct}
\end{figure}

\itindbeg{scheduler}

Una parte del kernel, lo \textit{scheduler}, si occupa di stabilire, sulla
base di un opportuno calcolo delle priorità e con una suddivisione appropriata
del tempo di processore, quali fra i vari ``\textsl{processi}'' presenti nel
sistema deve essere eseguito, realizzando il cosiddetto
\itindex{preemptive~multitasking} \textit{preemptive
  multitasking}.\footnote{si chiama così quella gestione del
  \textit{multitasking} in cui è il kernel a decidere a chi assegnare l'uso
  della CPU, potendo interrompere l'esecuzione di un processo in qualunque
  momento.}  Ogni processo verrà comunque eseguito in modalità protetta;
quando necessario esso potrà accedere alle risorse della macchina soltanto
attraverso delle ``\textsl{chiamate al sistema}'' (vedi
sez.~\ref{sec:intro_syscall}) che restituiranno il controllo al kernel per
eseguire le operazioni necessarie.

\itindend{scheduler}

La memoria viene sempre gestita dal kernel attraverso il meccanismo della
\textsl{memoria virtuale}, che consente di assegnare a ciascun processo uno
spazio di indirizzi ``\textsl{virtuale}'' (vedi sez.~\ref{sec:proc_memory})
che il kernel stesso, con l'ausilio della unità di gestione della memoria, si
incaricherà di rimappare automaticamente sulla memoria fisica disponibile, con
la possibilità ulteriore di spostare temporaneamente su disco (nella
cosiddetta area di \textit{swap}) parte di detta memoria qualora ci si trovi
nella necessità di liberare risorse.

Le periferiche infine vengono normalmente viste attraverso un'interfaccia
astratta che permette di trattarle come se fossero dei file, secondo uno dei
concetti base della architettura di Unix, per cui ``\textsl{tutto è in file}''
(\textit{everything is a file}) su cui torneremo in
sez.~\ref{sec:intro_file_dir}. In realtà questo non è sempre vero (ad esempio
non lo è per le interfacce di rete) dato che ci sono periferiche che non
rispondendo bene a questa astrazione richiedono un'interfaccia diversa.  Anche
in questo caso però resta valido il concetto generale che tutto il lavoro di
accesso e gestione delle periferiche a basso livello viene effettuato dal
kernel tramite l'opportuno codice di gestione delle stesse, che in
fig.~\ref{fig:intro_sys_struct} si è indicato come \textit{driver}.


\subsection{Il kernel e il sistema}
\label{sec:intro_kern_and_sys}

Uno dei concetti fondamentali su cui si basa l'architettura dei sistemi Unix è
quello della distinzione fra il cosiddetto \textit{user space}, che
contraddistingue l'ambiente in cui vengono eseguiti i programmi, e il
\textit{kernel space}, che è l'ambiente in cui viene eseguito il kernel. Ogni
programma vede sé stesso come se avesse la piena disponibilità della CPU e
della memoria ed è, salvo i meccanismi di comunicazione previsti dal sistema,
completamente ignaro del fatto che altri programmi possono essere messi in
esecuzione dal kernel.

Per questa separazione non è possibile ad un singolo programma disturbare
l'azione di un altro programma o del kernel stesso, e questo è il principale
motivo della stabilità di un sistema unix-like nei confronti di altri sistemi
in cui i processi non hanno di questi limiti o in cui essi vengono eseguiti
allo stesso livello del kernel. Pertanto deve essere chiaro a chi programma in
un sistema unix-like che l'accesso diretto all'hardware non può avvenire se
non all'interno del kernel; al di fuori dal kernel il programmatore deve usare
le opportune interfacce che quest'ultimo fornisce per i programmi in
\textit{user space}.

Per capire meglio la distinzione fra \textit{kernel space} e \textit{user
  space} si può prendere in esame la procedura di avvio di un sistema
unix-like. All'accensione il \textit{firmware} presente nella EPROM della
propria macchina (per i PC compatibili il BIOS), eseguirà la procedura di
avvio del sistema, il cosiddetto \textit{bootstrap},\footnote{il nome deriva
  da un'espressione gergale che significa ``sollevarsi da terra tirandosi per
  le stringhe delle scarpe'', per indicare il compito, almeno apparentemente
  impossibile, di far eseguire un programma a partire da un computer appena
  acceso che appunto non ne contiene nessuno; non è impossibile perché in
  realtà c'è un programma iniziale, che è il BIOS.} incaricandosi di caricare
il kernel in memoria e di farne partire l'esecuzione. 

A questo punto il controllo passerà al kernel, il quale però da parte sua, una
volta inizializzato opportunamente l'hardware, si limiterà a due sole
operazioni, montare il filesystem radice (torneremo su questo in
sez.~\ref{sec:file_arch_overview}) e lanciare il primo processo che eseguirà
il programma di inizializzazione del sistema, che in genere, visto il suo
scopo, si chiama \cmd{init}.

Una volta lanciato \cmd{init} tutto il lavoro successivo verrà eseguito
\textit{user space} da questo programma, che sua volta si incaricherà di
lanciare tutti gli altri programmi, fra cui ci sarà quello che si occupa di
dialogare con la tastiera e lo schermo della console, quello che mette a
disposizione un terminale e la \textit{shell} da cui inviare i comandi
all'utente che si vuole collegare, ed in generale tutto quanto necessario ad
avere un sistema utilizzabile.

E' da rimarcare come tutto ciò che riguarda l'interazione con l'utente, che
usualmente viene visto come parte del sistema, non abbia in realtà niente a
che fare con il kernel, ma sia effettuato da opportuni programmi che vengono
eseguiti, allo stesso modo di un qualunque programma di scrittura o di disegno
e della stessa interfaccia grafica, in \textit{user space}.

Questo significa ad esempio che il sistema di per sé non dispone di primitive
per tutta una serie di operazioni (ad esempio come la copia di un file) che
altri sistemi (come Windows) hanno invece al loro interno. Questo perché tutte
le operazioni di normale amministrazione di un sistema, sono effettuata
attraverso dei normali programmi utilizzando le interfacce di programmazione
che il kernel mette a disposizione.

È per questo motivo che quando ci si riferisce al sistema nella sua interezza
viene spesso sottolineato come sia corretto parlare di ``GNU/Linux'' e non di
Linux; da solo infatti il kernel non è sufficiente, quello che costruisce un
sistema operativo utilizzabile è la presenza di tutta una serie di librerie e
programmi di utilità, ed i più comuni sono appunto quelli realizzati dal
progetto GNU della Free Software Foundation, grazie ai quali si possono
eseguire le normali operazioni che ci si aspetta da un sistema operativo.


\subsection{\textit{System call} e funzioni di libreria}
\label{sec:intro_syscall}

Come illustrato in fig.~\ref{fig:intro_sys_struct} i programmi possono
accedere ai servizi forniti dal kernel tramite opportune interfacce dette
\textit{system call} (\textsl{chiamate al sistema}, appunto). Si tratta di un
insieme di funzioni che un programma può invocare, per le quali viene generata
un'interruzione nell'esecuzione del codice del processo, passando il controllo
al kernel. Sarà quest'ultimo che eseguirà in le operazioni relative alla
funzione richiesta in \textit{kernel space}, restituendo poi i risultati al
chiamante.

Ogni versione di Unix ha storicamente sempre avuto un certo numero di
\textit{system call}, che sono documentate nella seconda sezione del
\textsl{Manuale di programmazione di Unix}, quella cui si accede con il
comando \cmd{man 2 <nome>}, ed anche Linux non fa eccezione. Queste
\textit{system call} sono poi state codificate da vari standard, che
esamineremo brevemente in sez.~\ref{sec:intro_standard}.

Normalmente ciascuna \textit{system call} fornita dal kernel viene associata
ad una funzione con lo stesso nome definita all'interno della libreria
fondamentale del sistema, quella che viene chiamata \textsl{Libreria Standard
  del C} (\textit{C Standard Library}) in ragione del fatto che il primo
kernel Unix e tutti i programmi eseguiti su di esso vennero scritti in questo
linguaggio, usando le sue librerie. In seguito faremo riferimento alle
funzioni di questa libreria che si interfacciano alle \textit{system call}
come ``\textsl{funzioni di sistema}''.

Questa libreria infatti, oltre alle chiamate alle \textit{system call},
contiene anche tutta una serie di ulteriori funzioni di utilità che vengono
comunemente usate nella programmazione e sono definite nei vari standard che
documentano le interfacce di programmazione di un sistema unix-like. Questo
concetto è importante da tener presente perché programmare in Linux significa
anche essere in grado di usare le funzioni fornite dalla libreria Standard del
C, in quanto né il kernel né il linguaggio C implementano direttamente
operazioni ordinarie come l'allocazione dinamica della memoria, l'input/output
bufferizzato sui file, la manipolazione delle stringhe, la matematica in
virgola mobile, che sono comunemente usate da qualunque programma.

Tutto ciò mette nuovamente in evidenza il fatto che nella stragrande
maggioranza dei casi si dovrebbe usare il nome GNU/Linux in quanto una parte
essenziale del sistema, senza la quale niente funzionerebbe, è appunto la
\textit{GNU Standard C Library} (a cui faremo da qui in avanti riferimento
come \acr{glibc}), la libreria standard del C realizzata dalla Free Software
Foundation nella quale sono state implementate tutte le funzioni essenziali
definite negli standard POSIX e ANSI C (e molte altre), che vengono utilizzate
da qualunque programma.

Si tenga comunque presente che questo non è sempre vero, dato che esistono
implementazioni alternative della libreria standard del C, come la
\textit{libc5} o la \textit{uClib}, che non derivano dal progetto GNU. La
\textit{libc5}, che era usata con le prime versioni del kernel Linux, è oggi
ormai completamente soppiantata dalla \acr{glibc}. La \textit{uClib} invece,
pur non essendo completa come la \acr{glibc}, resta molto diffusa nel mondo
dei dispositivi \textit{embedded} per le sue dimensioni estremamente ridotte,
e soprattutto per la possibilità di togliere le parti non necessarie. Pertanto
costituisce un valido rimpiazzo della \acr{glibc} in tutti quei sistemi
specializzati che richiedono una minima occupazione di memoria. Infine per lo
sviluppo del sistema Android è stata realizzata da Google un'altra libreria
standard del C, utilizzata principalmente per evitare l'uso della \acr{glibc}.

Tradizionalmente le funzioni specifiche della libreria standard del C sono
riportate nella terza sezione del \textsl{Manuale di Programmazione di Unix}
(cioè accessibili con il comando \cmd{man 3 <nome>}) e come accennato non sono
direttamente associate ad una \textit{system call} anche se, ad esempio per la
gestione dei file o della allocazione dinamica della memoria, possono farne
uso nella loro implementazione.  Nonostante questa distinzione, fondamentale
per capire il funzionamento del sistema, l'uso da parte dei programmi di una
di queste funzioni resta lo stesso, sia che si tratti di una funzione interna
della libreria che di una \textit{system call}.


\subsection{Un sistema multiutente}
\label{sec:intro_multiuser}

Linux, come gli altri kernel Unix, nasce fin dall'inizio come sistema
multiutente, cioè in grado di fare lavorare più persone in contemporanea. Per
questo esistono una serie di meccanismi di sicurezza, che non sono previsti in
sistemi operativi monoutente, e che occorre tenere presenti. In questa sezione
parleremo brevemente soltanto dei meccanismi di sicurezza tradizionali di un
sistema unix-like, oggi molti di questi sono stati notevolmente estesi
rispetto al modello tradizionale, ma per il momento ignoreremo queste
estensioni.

Il concetto base è quello di utente (\textit{user}) del sistema, le cui
capacità sono sottoposte a limiti precisi.  Sono così previsti una serie di
meccanismi per identificare i singoli utenti ed una serie di permessi e
protezioni per impedire che utenti diversi possano danneggiarsi a vicenda o
danneggiare il sistema. Questi meccanismi sono realizzati dal kernel stesso ed
attengono alle operazioni più varie, e torneremo su di essi in dettaglio più
avanti.

Normalmente l'utente è identificato da un nome (il cosiddetto
\textit{username}), che ad esempio è quello che viene richiesto all'ingresso
nel sistema dalla procedura di \textit{login} (torneremo su questo in
sez.~\ref{sec:sess_login}).  Questa procedura si incarica di verificare
l'identità dell'utente, in genere attraverso la richiesta di una parola
d'ordine (la \textit{password}), anche se sono possibili meccanismi
diversi.\footnote{ad esempio usando la libreria PAM (\textit{Pluggable
    Autentication Methods}) è possibile generalizzare i meccanismi di
  autenticazione e sostituire ad esempio l'uso delle password con meccanismi
  di identificazione biometrica; per un approfondimento dell'argomento si
  rimanda alla sez.~4.3.7 di \cite{AGL}.} Eseguita la procedura di
riconoscimento in genere il sistema manda in esecuzione un programma di
interfaccia (che può essere la \textit{shell} su terminale o un'interfaccia
grafica) che mette a disposizione dell'utente un meccanismo con cui questo può
impartire comandi o eseguire altri programmi.

Ogni utente appartiene anche ad almeno un gruppo (il cosiddetto
\textit{default group}), ma può essere associato ad altri gruppi (i
\textit{supplementary group}), questo permette di gestire i permessi di
accesso ai file e quindi anche alle periferiche, in maniera più flessibile,
definendo gruppi di lavoro, di accesso a determinate risorse, ecc. 

L'utente e il gruppo sono identificati dal kernel un identificativo numerico,
la cui corrispondenza ad un nome espresso in caratteri è inserita nei due file
\conffile{/etc/passwd} e \conffile{/etc/group}.\footnote{in realtà negli
  sistemi più moderni, come vedremo in sez.~\ref{sec:sys_user_group} queste
  informazioni possono essere mantenute, con l'uso del \textit{Name Service
    Switch}, su varie tipologie di supporti, compresi server centralizzati
  come LDAP.}  Questi identificativi sono l'\textit{user identifier}, detto in
breve \textsl{user-ID}, ed indicato dall'acronimo \ids{UID}, e il
\textit{group identifier}, detto in breve \textsl{group-ID}, ed identificato
dall'acronimo \ids{GID}, torneremo in dettaglio su questo argomento in
sez.~\ref{sec:proc_perms}.  Il kernel conosce ed utilizza soltanto questi
valori numerici, i nomi ad essi associati sono interamente gestiti in
\textit{user space} con opportune funzioni di libreria, torneremo su questo
argomento in sez.~\ref{sec:sys_user_group}.
 
Grazie a questi identificativi il sistema è in grado di tenere traccia
dell'utente a cui appartiene ciascun processo ed impedire ad altri utenti di
interferire con quest'ultimo.  Inoltre con questo sistema viene anche
garantita una forma base di sicurezza interna in quanto anche l'accesso ai
file (vedi sez.~\ref{sec:file_access_control}) è regolato da questo meccanismo
di identificazione.

Infine in ogni sistema unix-like è presente uno speciale utente privilegiato,
il cosiddetto \textit{superuser}, il cui username è di norma \textit{root}, ed
il cui \ids{UID} è zero. Esso identifica l'amministratore del sistema, che
deve essere in grado di fare qualunque operazione; per l'utente \textit{root}
infatti i meccanismi di controllo cui si è accennato in precedenza sono
disattivati.\footnote{i controlli infatti vengono eseguiti da uno pseudo-codice
  del tipo: ``\code{if (uid) \{ \textellipsis\ \}}''.}


\section{L'architettura di file e directory}
\label{sec:intro_file_dir}

Come accennato in sez.~\ref{sec:intro_base_concept} uno dei concetti
fondamentali dell'architettura di un sistema Unix è il cosiddetto
\textit{everything is a file} (\textsl{tutto è un file}), cioè il fatto che
l'accesso ai vari dispositivi di input/output del computer viene effettuato
attraverso un'interfaccia astratta che tratta le periferiche allo stesso modo
dei normali file di dati.

In questa sezione forniremo una descrizione a grandi linee dell'architettura
della gestione dei file in Linux, partendo da una introduzione ai concetti di
base, per poi illustrare la struttura dell'albero dei file ed il significato
dei tipi di file, concludendo con una panoramica sulle caratteristiche
principali delle due interfacce con cui i processi possono effettuare l'I/O su
file.


\subsection{Una panoramica generale}
\label{sec:file_arch_overview}

Per poter accedere ai file, il kernel deve mettere a disposizione dei
programmi delle opportune \textit{system call} che consentano di leggere e
scrivere il contenuto. Tutto ciò ha due aspetti: il primo è che il kernel, per
il concetto dell'\textit{everything is a file}, deve fornire una interfaccia
che consenta di operare sui file, sia che questi corrispondano ai normali file
di dati, o ai cosiddetti ``\textsl{file speciali}'', come i file di
dispositivo (o \textit{device file}) che permettono di accedere alle
periferiche o le \textit{fifo} ed i socket che forniscono funzionalità di
comunicazione fra processi (torneremo su questo in sez.~\ref{sec:file_mknod}).

Il secondo aspetto è che per poter utilizzare dei normali file di dati il
kernel deve provvedere ad organizzare e rendere accessibile in maniera
opportuna l'informazione in essi contenuta memorizzandola sullo spazio grezzo
disponibile sui dischi.  Questo viene fatto strutturando l'informazione sul
disco attraverso quello che si chiama un
``\textit{filesystem}''. L'informazione così strutturata poi viene resa
disponibile ai processi attraverso quello che viene chiamato il
``\textsl{montaggio}'' del filesystem nell'albero dei file, dove il contenuto
sarà accessibile nella forma ordinaria di file e directory.

\itindbeg{Virtual~File~System~(VFS)}

In Linux il concetto di \textit{everything is a file} è stato implementato
attraverso il cosiddetto \textit{Virtual File System} (da qui in avanti VFS)
che è uno strato intermedio che il kernel usa per accedere ai più svariati
filesystem mantenendo la stessa interfaccia per i programmi in \textit{user
  space}.  Il VFS fornisce quel livello di astrazione che permette di
collegare le operazioni interne del kernel per la manipolazione sui file con
le \textit{system call} relative alle operazioni di I/O, e gestisce poi
l'organizzazione di dette operazioni nei vari modi in cui i diversi filesystem
le effettuano, permettendo la coesistenza di filesystem differenti all'interno
dello stesso albero delle directory. Approfondiremo il funzionamento di
interfaccia generica fornita dal VFS in sez.~\ref{sec:file_vfs_work}.

In sostanza quello che accade è che quando un processo esegue una
\textit{system call} che opera su un file, il kernel chiama sempre una
funzione implementata nel VFS. La funzione eseguirà le manipolazioni sulle
strutture generiche e utilizzerà poi delle chiamate alle opportune funzioni
del filesystem specifico a cui si fa riferimento. Saranno queste infine a
chiamare le funzioni di più basso livello che eseguono le operazioni di I/O
sul dispositivo fisico, secondo lo schema riportato in
fig.~\ref{fig:file_VFS_scheme}.

\begin{figure}[!htb]
  \centering
  \includegraphics[width=8cm]{img/vfs}
  \caption{Schema delle operazioni del VFS.}
  \label{fig:file_VFS_scheme}
\end{figure}

Questa interfaccia resta la stessa anche quando, invece che a dei normali
file, si accede alle periferiche coi citati file di dispositivo, solo che in
questo caso invece di usare il codice del filesystem che accede al disco, il
\textit{Virtual File System} eseguirà direttamente il codice del kernel che
permette di accedere alla periferica.

\itindend{Virtual~File~System~(VFS)}

Come accennato in precedenza una delle funzioni essenziali per il
funzionamento dell'interfaccia dei file è quella che consente di montare un
filesystem nell'albero dei file, e rendere così visibili i suoi contenuti. In
un sistema unix-like infatti, a differenza di quanto avviene in altri sistemi
operativi, tutti i file vengono mantenuti all'interno di un unico albero la
cui radice (quella che viene chiamata \textit{root directory}) viene montata
all'avvio direttamente dal kernel.

Come accennato in sez.~\ref{sec:intro_kern_and_sys}) montare la radice è,
insieme al lancio di \cmd{init},\footnote{l'operazione è ovviamente anche
  preliminare al lancio di \cmd{init}, dato il kernel deve poter accedere al
  file che contiene detto programma.} l'unica operazione che viene effettuata
direttamente dal kernel in fase di avvio. Il kernel infatti, completata la
fase di inizializzazione, utilizza l'indicazione passatagli dal
\textit{bootloader} su quale sia il dispositivo che contiene il filesystem da
usare come punto di partenza, lo monta come radice dell'albero dei file.

Tutti gli ulteriori filesystem che possono essere disponibili su altri
dispositivi dovranno a loro volta essere inseriti nell'albero, montandoli in
altrettante directory del filesystem radice, su quelli che vengono chiamati
\textit{mount point}.  Questo comunque avverrà sempre in un secondo tempo, a
cura dei programmi eseguiti nella procedura di inizializzazione del sistema,
grazie alle funzioni che tratteremo in sez.~\ref{sec:filesystem_mounting}.


\subsection{La risoluzione del nome di file e directory}
\label{sec:file_pathname}

\itindbeg{pathname}

Come appena illustrato sez.~\ref{sec:file_arch_overview} una delle
caratteristiche distintive di un sistema unix-like è quella di avere un unico
albero dei file. Un file deve essere identificato dall'utente usando quello
che viene chiamato il suo \textit{pathname},\footnote{il manuale della
  \acr{glibc} depreca questa nomenclatura, che genererebbe confusione poiché
  \textit{path} indica anche un insieme di directory su cui effettuare una
  ricerca (come quello in cui la shell cerca i comandi). Al suo posto viene
  proposto l'uso di \textit{filename} e di componente per il nome del file
  all'interno della directory. Non seguiremo questa scelta dato che l'uso
  della parola \textit{pathname} è ormai così comune che mantenerne l'uso è
  senz'altro più chiaro dell'alternativa proposta.} vale a dire tramite il
``\textsl{percorso}'' (nome che talvolta viene usato come traduzione di
\textit{pathname}) che si deve fare per accedere al file a partire da una
certa ``\textit{directory}''.

Una directory in realtà è anch'essa un file, nel senso che è anch'essa un
oggetto di un filesystem, solo che è un file particolare che il kernel
riconosce appositamente come tale per poterlo utilizzare come directory. Il
suo scopo è quello di contenere una lista di nomi di file e le informazioni
che associano ciascuno di questi nomi al relativo contenuto (torneremo su
questo in sez.~\ref{sec:file_arch_func}).

Dato che questi nomi possono corrispondere ad un qualunque altro oggetto del
filesystem, compresa un'altra directory, si ottiene naturalmente
un'organizzazione ad albero inserendo nomi di directory dentro altre
directory.  All'interno dello stesso albero si potranno poi inserire anche
tutti gli altri oggetti previsti l'interfaccia del VFS (su cui torneremo in
sez.~\ref{sec:file_file_types}), come le \textit{fifo}, i collegamenti
simbolici, i socket e gli stessi file di dispositivo.

La convenzione usata nei sistemi unix-like per indicare i \textit{pathname}
dei file è quella di usare il carattere ``\texttt{/}'' come separatore fra i
nomi che indicano le directory che lo compongono. Dato che la directory radice
sta in cima all'albero, essa viene indicata semplicemente con il
\textit{pathname} ``\file{/}''.

\itindbeg{pathname~resolution}

Un file può essere indicato rispetto ad una directory semplicemente
specificandone il nome, il manuale della \acr{glibc} chiama i nomi contenuti
nelle directory ``componenti'' (in inglese \textit{file name components}), noi
li chiameremo più semplicemente \textsl{nomi} o \textsl{voci}, riservando la
parola \textsl{componenti} ai nomi che, separati da una ``\texttt{/}'',
costituiscono il \textit{pathname}. Questi poi dovranno corrispondere, perché
il \textit{pathname} sia valido, a voci effettivamente presenti nelle
directory, ma non è detto che un \textit{pathname} debba per forza risultare
valido.  

Il procedimento con cui dato un \textit{pathname} si individua il file a cui
esso fa riferimento, è chiamato \textsl{risoluzione del percorso}
(\textit{filename resolution} o \textit{pathname resolution}). Lo stesso
procedimento ci può anche dire che il \textit{pathname} usato non è valido.
La risoluzione viene eseguita esaminando il \textit{pathname} da sinistra a
destra e localizzando ogni componente dello stesso come nome in una directory
a partire dalla directory iniziale, usando il carattere ``\texttt{/}'' come
separatore per scendere dall'una all'altra. Nel caso si indichi un componente
vuoto il costrutto ``\texttt{//}'' viene considerato equivalente a
``\texttt{/}''.

Ovviamente perché la risoluzione abbia successo occorre che i componenti
intermedi esistano e siano effettivamente directory, e che il file o la
directory indicata dall'ultimo componente esista.  Inoltre i permessi relativi
alle directory indicate nel \textit{pathname} (torneremo su questo
sez.~\ref{sec:file_access_control}) dovranno consentire l'accesso all'intero
\textit{pathname}.

\itindsubbeg{pathname}{assoluto}
\itindsubbeg{pathname}{relativo}

Se il \textit{pathname} comincia con il carattere ``\texttt{/}'' la ricerca
parte dalla directory radice del processo. Questa, a meno di non avere
eseguito una \func{chroot} (funzione su cui torneremo in
sez.~\ref{sec:file_chroot}) è la stessa per tutti i processi ed equivale alla
directory radice dell'albero dei file montata dal kernel all'avvio del
sistema; in questo caso si parla di un \textsl{pathname assoluto}. Altrimenti
la ricerca parte dalla directory di lavoro corrente del processo (su cui
torneremo in sez.~\ref{sec:file_work_dir}) ed il \textit{pathname} è detto
\textsl{pathname relativo}.

\itindsubend{pathname}{assoluto}
\itindsubend{pathname}{relativo}

Infine i nomi di directory ``\file{.}'' e ``\file{..}'' hanno un significato
speciale e vengono inseriti in ogni directory quando questa viene creata (vedi
sez.~\ref{sec:file_dir_creat_rem}). Il primo fa riferimento alla directory
corrente e il secondo alla directory \textsl{genitrice} (o \textit{parent
  directory}) cioè la directory che contiene il riferimento alla directory
corrente.

In questo modo con ``\file{..}'' si può usare un \textit{pathname} relativo
per indicare un file posto al di sopra della directory corrente, tornando
all'indietro nell'albero dei file.  Questa retromarcia però su fermerà una
volta raggiunta la directory radice, perché non esistendo in questo caso una
directory superiore, il nome ``\file{..}''  farà riferimento alla radice
stessa.

\itindend{pathname}
\itindend{pathname~resolution}


\subsection{I tipi di file}
\label{sec:file_file_types}

Parlare dei tipi di file su Linux, come per qualunque sistema unix-like,
significa anzitutto chiarire il proprio vocabolario e sottolineare le
differenze che ci sono rispetto ad altri sistemi operativi.

\index{file!di~dispositivo|(}
\index{file!speciali|(} 

Come accennato in sez.~\ref{sec:file_arch_overview} su Linux l'uso del
\textit{Virtual File System} consente di trattare come file oggetti molto
diversi fra loro. Oltre ai normali file di dati abbiamo già accennato ad altri
due di questi oggetti, i file di dispositivo e le directory, ma ne esistono
altri. In genere quando si parla di tipo di file su Linux si fa riferimento a
questi diversi tipi, di cui si riportato l'elenco completo in
tab.~\ref{tab:file_file_types}.

\begin{table}[htb]
  \footnotesize
  \centering
    \begin{tabular}[c]{|l|l|p{6cm}|}
    \hline
    \multicolumn{2}{|c|}{\textbf{Tipo di file}} & \textbf{Descrizione} \\
    \hline
    \hline
      \textit{regular file}& \textsl{file regolare}
                           & Un file che contiene dei dati (l'accezione
                             normale di file).\\
      \textit{directory}   &\textsl{cartella o direttorio}
                           & Un file che contiene una lista di nomi associati
                             a degli \textit{inode} (vedi
                             sez.~\ref{sec:file_vfs_work}).\\   
      \textit{symbolic link}&\textsl{collegamento simbolico}
                           & Un file che contiene un riferimento ad un altro 
                             file/directory.\\ 
      \textit{char device} &\textsl{dispositivo a caratteri} 
                           & Un file \textsl{speciale} che identifica una
                             periferica ad accesso a caratteri.\\
      \textit{block device}& \textsl{dispositivo a blocchi} 
                           & Un file \textsl{speciale} che identifica una
                             periferica ad accesso a blocchi.\\
      \textit{fifo}        & ``\textsl{coda}'' 
                           & Un file \textsl{speciale} che identifica una
                             linea di comunicazione unidirezionale (vedi
                             sez.~\ref{sec:ipc_named_pipe}).\\
      \textit{socket}      & ``\textsl{presa}''
                           & Un file \textsl{speciale} che identifica una
                             linea di comunicazione bidirezionale (vedi
                             cap.~\ref{cha:socket_intro}).\\
    \hline
    \end{tabular}
    \caption{Tipologia dei file definiti nel VFS}
    \label{tab:file_file_types}
\end{table}

Si tenga ben presente che questa classificazione non ha nulla a che fare con
una classificazione dei file in base al loro contenuto, dato che in tal caso
si avrebbe a che fare sempre e solo con dei file di dati. E non ha niente a
che fare neanche con le eventuali diverse modalità con cui si possa accedere
al contenuto dei file di dati.  La classificazione di
tab.~\ref{tab:file_file_types} riguarda il tipo di oggetti gestiti dal
\textit{Virtual File System}, ed è da notare la presenza dei cosiddetti file
``\textsl{speciali}''.

Alcuni di essi, come le \textit{fifo} (che tratteremo in
sez.~\ref{sec:ipc_named_pipe}) ed i \textit{socket} (che tratteremo in
cap.~\ref{cha:socket_intro}) non sono altro che dei riferimenti per utilizzare
alcune funzionalità di comunicazione fornite dal kernel. Gli altri sono
proprio quei \textsl{file di dispositivo} che costituiscono una interfaccia
diretta per leggere e scrivere sui dispositivi fisici. Anche se finora li
abbiamo chiamati genericamente così, essi sono tradizionalmente suddivisi in
due grandi categorie, \textsl{a blocchi} e \textsl{a caratteri}, a seconda
delle modalità in cui il dispositivo sottostante effettua le operazioni di
I/O.

I dispositivi a blocchi (ad esempio i dischi) sono quelli corrispondono a
periferiche per le quali è richiesto che l'I/O venga effettuato per blocchi di
dati di dimensioni fissate (nel caso dei dischi le dimensioni di un settore),
mentre i dispositivi a caratteri sono quelli per cui l'I/O può essere
effettuato senza nessuna particolare struttura, ed in generale anche un byte
alla volta, da cui il nome.

Una delle differenze principali con altri sistemi operativi come il VMS o
Windows è che per Unix tutti i file di dati sono identici e contengono un
flusso continuo di byte. Non esiste cioè differenza per come vengono visti dal
sistema file di diverso contenuto o formato, come nel caso di quella fra file
di testo e binari che c'è in Windows. Non c'è neanche una strutturazione a
record per il cosiddetto ``\textsl{accesso diretto}'' come nel caso del
VMS.\footnote{questo vale anche per i dispositivi a blocchi: la strutturazione
  dell'I/O in blocchi di dimensione fissa avviene solo all'interno del kernel,
  ed è completamente trasparente all'utente; inoltre talvolta si parla di
  \textsl{accesso diretto} riferendosi alla capacità, che non ha niente a che
  fare con tutto ciò, di effettuare attraverso degli appositi file di
  dispositivo delle operazioni di I/O direttamente sui dischi senza passare
  attraverso un filesystem, il cosiddetto \textit{raw access}, introdotto coi
  kernel della serie 2.4.x ma ormai in sostanziale disuso.}

\index{file!di~dispositivo|)}
\index{file!speciali|)} 

Una differenza che attiene ai contenuti di un file però esiste, ed è relativa
al formato dei file di testo. Nei sistemi unix-like la fine riga è codificata
in maniera diversa da Windows o dal vecchio MacOS, in particolare il fine riga
è il carattere \texttt{LF} (\verb|\n|) al posto del \texttt{CR} (\verb|\r|)
del vecchio MacOS e del \texttt{CR LF} (\verb|\r\n|) di Windows. Questo può
causare alcuni problemi qualora nei programmi si facciano assunzioni sul
terminatore della riga e per questo esistono dei programmi come \cmd{unix2dos}
e \cmd{dos2unix} che effettuano una conversione fra questi due formati di
testo.

Si ricordi comunque che un kernel unix-like non fornisce nessun supporto per
la tipizzazione dei file di dati in base al loro contenuto e che non c'è
nessun supporto per una qualche interpretazione delle estensioni (nel nome del
file) da parte del kernel,\footnote{non è così ad esempio nel filesystem HFS
  dei Mac, che supporta delle risorse associate ad ogni file, che specificano
  fra l'altro il contenuto ed il programma da usare per leggerlo; in realtà
  per alcuni filesystem esiste la possibilità di associare delle risorse ai
  file con gli \textit{extended attributes} (vedi sez.~\ref{sec:file_xattr}),
  ma è una caratteristica tutt'ora poco utilizzata, dato che non corrisponde
  al modello classico dei file in un sistema Unix.} ogni classificazione di
questo tipo avviene sempre in \textit{user space}. Gli unici file di cui il
kernel deve essere in grado di capire il contenuto sono i binari dei
programmi, per i quali sono supportati solo alcuni formati, anche se oggi
viene usato quasi esclusivamente
\itindex{Executable~and~Linkable~Format~(ELF)} l'ELF.\footnote{il nome è
  l'acronimo di \textit{Executable and Linkable Format}, un formato per
  eseguibili binari molto flessibile ed estendibile definito nel 1995 dal
  \textit{Tool Interface Standard} che per le sue caratteristiche di non
  essere legato a nessun tipo di processore o architettura è stato adottato da
  molti sistemi unix-like e non solo.}

\itindbeg{magic~number}

Nonostante l'assenza di supporto da parte del kernel per la classificazione
del contenuto dei file di dati, molti programmi adottano comunque delle
convenzioni per i nomi dei file, ad esempio il codice C normalmente si mette
in file con l'estensione ``\file{.c}''. Inoltre una tecnica molto usata per
classificare i contenuti da parte dei programmi è quella di utilizzare i primi
byte del file per memorizzare un ``\textit{magic number}'' che ne classifichi
il contenuto. Il concetto è quello di un numero intero, solitamente fra 2 e 10
byte, che identifichi il contenuto seguente, dato che questi sono anche
caratteri è comune trovare espresso tale numero con stringhe come
``\texttt{\%PDF}'' per i PDF o ``\texttt{\#!}'' per gli script. Entrambe
queste tecniche, per quanto usate ed accettate in maniera diffusa, restano
solo delle convenzioni il cui rispetto è demandato alle applicazioni stesse.

\itindend{magic~number}


\subsection{Le due interfacce per l'accesso ai file}
\label{sec:file_io_api}


\itindbeg{file~descriptor}

In Linux le interfacce di programmazione per l'I/O su file due.  La prima è
l'interfaccia nativa del sistema, quella che il manuale della \textsl{glibc}
chiama interfaccia dei ``\textit{file descriptor}'' (in italiano
\textsl{descrittori di file}). Si tratta di un'interfaccia specifica dei
sistemi unix-like che fornisce un accesso non bufferizzato.

L'interfaccia è essenziale, l'accesso viene detto non bufferizzato in quanto
la lettura e la scrittura vengono eseguite chiamando direttamente le
\textit{system call} del kernel, anche se in realtà il kernel effettua al suo
interno alcune bufferizzazioni per aumentare l'efficienza nell'accesso ai
dispositivi. L'accesso viene gestito attraverso i \textit{file descriptor} che
sono rappresentati da numeri interi (cioè semplici variabili di tipo
\ctyp{int}).  L'interfaccia è definita nell'\textit{header file}
\headfile{unistd.h} e la tratteremo in dettaglio in
sez.~\ref{sec:file_unix_interface}.

\itindbeg{file~stream}

La seconda interfaccia è quella che il manuale della \acr{glibc} chiama dei
\textit{file stream} o più semplicemente degli \textit{stream}.\footnote{in
  realtà una interfaccia con lo stesso nome è stata introdotta a livello di
  kernel negli Unix derivati da \textit{System V}, come strato di astrazione
  per file e socket; in Linux questa interfaccia, che comunque ha avuto poco
  successo, non esiste, per cui facendo riferimento agli \textit{stream}
  useremo il significato adottato dal manuale della \acr{glibc}.} Essa
fornisce funzioni più evolute e un accesso bufferizzato, controllato dalla
implementazione fatta nella \acr{glibc}.  Questa è l'interfaccia specificata
dallo standard ANSI C e perciò si trova anche su tutti i sistemi non Unix. Gli
\textit{stream} sono oggetti complessi e sono rappresentati da puntatori ad un
opportuna struttura definita dalle librerie del C, a cui si accede sempre in
maniera indiretta utilizzando il tipo \code{FILE *}; l'interfaccia è definita
nell'\textit{header file} \headfile{stdio.h} e la tratteremo in dettaglio in
sez.~\ref{sec:files_std_interface}.

Entrambe le interfacce possono essere usate per l'accesso ai file come agli
altri oggetti del VFS, ma per poter accedere alle operazioni di controllo
(descritte in sez.~\ref{sec:file_fcntl_ioctl}) su un qualunque tipo di oggetto
del VFS occorre usare l'interfaccia standard di Unix con i file
descriptor. Allo stesso modo devono essere usati i file descriptor se si vuole
ricorrere a modalità speciali di I/O come il \textit{file locking} o l'I/O
non-bloccante (vedi cap.~\ref{cha:file_advanced}).

Gli \textit{stream} forniscono un'interfaccia di alto livello costruita sopra
quella dei \textit{file descriptor}, che permette di poter scegliere tra
diversi stili di bufferizzazione.  Il maggior vantaggio degli \textit{stream}
è che l'interfaccia per le operazioni di input/output è molto più ricca di
quella dei \textit{file descriptor}, che forniscono solo funzioni elementari
per la lettura/scrittura diretta di blocchi di byte.  In particolare gli
\textit{stream} dispongono di tutte le funzioni di formattazione per l'input e
l'output adatte per manipolare anche i dati in forma di linee o singoli
caratteri.

In ogni caso, dato che gli \textit{stream} sono implementati sopra
l'interfaccia standard di Unix, è sempre possibile estrarre il \textit{file
  descriptor} da uno \textit{stream} ed eseguirvi sopra operazioni di basso
livello, o associare in un secondo tempo uno \textit{stream} ad un
\textit{file descriptor} per usare l'interfaccia più sofisticata.

In generale, se non necessitano specificatamente le funzionalità di basso
livello, è opportuno usare sempre gli \textit{stream} per la loro maggiore
portabilità, essendo questi ultimi definiti nello standard ANSI C;
l'interfaccia con i \textit{file descriptor} infatti segue solo lo standard
POSIX.1 dei sistemi Unix, ed è pertanto di portabilità più limitata.

\itindend{file~descriptor}
\itindend{file~stream}

\section{Gli standard}
\label{sec:intro_standard}

In questa sezione faremo una breve panoramica relativa ai vari standard che
nel tempo sono stati formalizzati da enti, associazioni, consorzi e
organizzazioni varie al riguardo ai sistemi operativi di tipo Unix o alle
caratteristiche che si sono stabilite come standard di fatto in quanto facenti
parte di alcune implementazioni molto diffuse come BSD o System V.

Ovviamente prenderemo in considerazione solo gli standard riguardanti
interfacce di programmazione e le altre caratteristiche di un sistema
unix-like (alcuni standardizzano pure i comandi base del sistema e la shell)
ed in particolare ci concentreremo sul come ed in che modo essi sono
supportati sia per quanto riguarda il kernel che la libreria standard del C,
con una particolare attenzione alla \acr{glibc}.


\subsection{Lo standard ANSI C}
\label{sec:intro_ansiC}

\itindbeg{ANSI~C}
Lo standard ANSI C è stato definito nel 1989 dall'\textit{American National
  Standard Institute} come prima standardizzazione del linguaggio C e per
questo si fa riferimento ad esso anche come C89. L'anno successivo è stato
adottato dalla ISO (\textit{International Standard Organisation}) come
standard internazionale con la sigla ISO/IEC 9899:1990, e per questo è noto
anche sotto il nome di standard ISO C, o ISO C90.  Nel 1999 è stata pubblicata
una revisione dello standard C89, che viene usualmente indicata come C99,
anche questa è stata ratificata dalla ISO con la sigla ISO/IEC 9899:1990, per
cui vi si fa riferimento anche come ISO C99.

Scopo dello standard è quello di garantire la portabilità dei programmi C fra
sistemi operativi diversi, ma oltre alla sintassi ed alla semantica del
linguaggio C (operatori, parole chiave, tipi di dati) lo standard prevede
anche una libreria di funzioni che devono poter essere implementate su
qualunque sistema operativo.

Per questo motivo, anche se lo standard non ha alcun riferimento ad un sistema
di tipo Unix, GNU/Linux (per essere precisi la \acr{glibc}), come molti Unix
moderni, provvede la compatibilità con questo standard, fornendo le funzioni
di libreria da esso previste. Queste sono dichiarate in una serie di
\textit{header file} anch'essi forniti dalla \acr{glibc} (tratteremo
l'argomento in sez.~\ref{sec:proc_syscall}).

In realtà la \acr{glibc} ed i relativi \textit{header file} definiscono un
insieme di funzionalità in cui sono incluse come sottoinsieme anche quelle
previste dallo standard ANSI C. È possibile ottenere una conformità stretta
allo standard (scartando le funzionalità addizionali) usando il \cmd{gcc} con
l'opzione \cmd{-ansi}. Questa opzione istruisce il compilatore a definire nei
vari \textit{header file} soltanto le funzionalità previste dallo standard
ANSI C e a non usare le varie estensioni al linguaggio e al preprocessore da
esso supportate.
\itindend{ANSI~C}


\subsection{I tipi di dati primitivi}
\label{sec:intro_data_types}

Uno dei problemi di portabilità del codice più comune è quello dei tipi di
dati utilizzati nei programmi, che spesso variano da sistema a sistema, o
anche da una architettura ad un'altra (ad esempio passando da macchine con
processori 32 bit a 64). In particolare questo è vero nell'uso dei cosiddetti
\index{tipo!elementare} \textit{tipi elementari} del linguaggio C (come
\ctyp{int}) la cui dimensione varia a seconda dell'architettura hardware.

Storicamente alcuni tipi nativi dello standard ANSI C sono sempre stati
associati ad alcune variabili nei sistemi Unix, dando per scontata la
dimensione. Ad esempio la posizione corrente all'interno di un file è stata
associata ad un intero a 32 bit, mentre il numero di dispositivo è stato
associato ad un intero a 16 bit. Storicamente questi erano definiti
rispettivamente come \ctyp{int} e \ctyp{short}, ma tutte le volte che, con
l'evolversi ed il mutare delle piattaforme hardware, alcuni di questi tipi si
sono rivelati inadeguati o sono cambiati, ci si è trovati di fronte ad una
infinita serie di problemi di portabilità.

\begin{table}[htb]
  \footnotesize
  \centering
  \begin{tabular}[c]{|l|l|}
    \hline
    \textbf{Tipo} & \textbf{Contenuto} \\
    \hline
    \hline
    \typed{caddr\_t} & Core address.\\
    \typed{clock\_t} & Contatore del \textit{process time} (vedi
                      sez.~\ref{sec:sys_cpu_times}.\\ 
    \typed{dev\_t}   & Numero di dispositivo (vedi sez.~\ref{sec:file_mknod}).\\
    \typed{gid\_t}   & Identificatore di un gruppo (vedi
                      sez.~\ref{sec:proc_access_id}).\\
    \typed{ino\_t}   & Numero di \textit{inode} 
                      (vedi sez.~\ref{sec:file_vfs_work}).\\ 
    \typed{key\_t}   & Chiave per il System V IPC (vedi
                      sez.~\ref{sec:ipc_sysv_generic}).\\
    \typed{loff\_t}  & Posizione corrente in un file.\\
    \typed{mode\_t}  & Attributi di un file.\\
    \typed{nlink\_t} & Contatore dei collegamenti su un file.\\
    \typed{off\_t}   & Posizione corrente in un file.\\
    \typed{pid\_t}   & Identificatore di un processo (vedi
                      sez.~\ref{sec:proc_pid}).\\
    \typed{rlim\_t}  & Limite sulle risorse.\\
    \typed{sigset\_t}& Insieme di segnali (vedi sez.~\ref{sec:sig_sigset}).\\
    \typed{size\_t}  & Dimensione di un oggetto.\\
    \typed{ssize\_t} & Dimensione in numero di byte ritornata dalle funzioni.\\
    \typed{ptrdiff\_t}& Differenza fra due puntatori.\\
    \typed{time\_t}  & Numero di secondi (in \textit{calendar time}, vedi 
                      sez.~\ref{sec:sys_time}).\\
    \typed{uid\_t}   & Identificatore di un utente (vedi
                      sez.~\ref{sec:proc_access_id}).\\
    \hline
  \end{tabular}
  \caption{Elenco dei tipi primitivi, definiti in \headfile{sys/types.h}.}
  \label{tab:intro_primitive_types}
\end{table}

Per questo motivo tutte le funzioni di libreria di solito non fanno
riferimento ai tipi elementari dello standard del linguaggio C, ma ad una
serie di \index{tipo!primitivo} \textsl{tipi primitivi} del sistema, riportati
in tab.~\ref{tab:intro_primitive_types}, e definiti nell'\textit{header file}
\headfiled{sys/types.h}, in modo da mantenere completamente indipendenti i tipi
utilizzati dalle funzioni di sistema dai tipi elementari supportati dal
compilatore C.


\subsection{Lo standard System V}
\label{sec:intro_sysv}

Come noto Unix nasce nei laboratori della AT\&T, che ne registrò il nome come
marchio depositato, sviluppandone una serie di versioni diverse; nel 1983 la
versione supportata ufficialmente venne rilasciata al pubblico con il nome di
Unix System V, e si fa rifermento a questa implementazione con la sigla SysV o
SV.

Negli anni successivi l'AT\&T proseguì lo sviluppo del sistema rilasciando
varie versioni con aggiunte e integrazioni, ed in particolare la
\textit{release 2} nel 1985, a cui si fa riferimento con il nome SVr2 e la
\textit{release 3} nel 1986 (denominata SVr3). Le interfacce di programmazione
di queste due versioni vennero descritte formalmente in due documenti
denominati \itindex{System~V~Interface~Definition~(SVID)} \textit{System V
  Interface Definition} (o SVID), pertanto nel 1985 venne rilasciata la
specifica SVID 1 e nel 1986 la specifica SVID 2.

Nel 1989 un accordo fra vari venditori (AT\&T, Sun, HP, ed altri) portò ad una
versione di System V che provvedeva un'unificazione delle interfacce
comprendente anche Xenix e BSD, questa venne denominata \textit{release 4} o
SVr4. Anche le relative interfacce vennero descritte in un documento dal
titolo \textit{System V Interface Description}, venendo a costituire lo
standard SVID 3, che viene considerato la specifica finale di System V, ed a
cui spesso si fa riferimento semplicemente con SVID. Anche SVID costituisce un
sovrainsieme delle interfacce definite dallo standard POSIX.  

Nel 1992 venne rilasciata una seconda versione del sistema, la SVr4.2; l'anno
successivo la divisione della AT\&T (già a suo tempo rinominata in Unix System
Laboratories) venne acquistata dalla Novell, che poi trasferì il marchio Unix
al consorzio X/Open. L'ultima versione di System V fu la SVr4.2MP rilasciata
nel dicembre 1993. Infine nel 1995 è stata rilasciata da SCO, che aveva
acquisito alcuni diritti sul codice di System V, una ulteriore versione delle
\textit{System V Interface Description}, che va sotto la denominazione di SVID
4.

Linux e la \acr{glibc} implementano le principali funzionalità richieste dalle
specifiche SVID che non sono già incluse negli standard POSIX ed ANSI C, per
compatibilità con lo Unix System V e con altri Unix (come SunOS) che le
includono. Tuttavia le funzionalità più oscure e meno utilizzate (che non sono
presenti neanche in System V) sono state tralasciate.

Le funzionalità implementate sono principalmente il meccanismo di
intercomunicazione fra i processi e la memoria condivisa (il cosiddetto System
V IPC, che vedremo in sez.~\ref{sec:ipc_sysv}) le funzioni della famiglia
\funcm{hsearch} e \funcm{drand48}, \funcm{fmtmsg} e svariate funzioni
matematiche.


\subsection{Lo ``\textsl{standard}'' BSD}
\label{sec:intro_bsd}

Lo sviluppo di BSD iniziò quando la fine della collaborazione fra l'Università
di Berkeley e la AT\&T generò una delle prime e più importanti fratture del
mondo Unix.  L'università di Berkeley proseguì nello sviluppo della base di
codice di cui disponeva, e che presentava parecchie migliorie rispetto alle
versioni allora disponibili, fino ad arrivare al rilascio di una versione
completa di Unix, chiamata appunto BSD, del tutto indipendente dal codice
della AT\&T.

Benché BSD non sia mai stato uno standard formalizzato, l'implementazione
dello Unix dell'Università di Berkeley nella sua storia ha introdotto una
serie di estensioni e interfacce di grandissima rilevanza, come i collegamenti
simbolici, la funzione \code{select} ed i socket di rete. Per questo motivo si
fa spesso riferimento esplicito alle interfacce presenti nelle varie versioni
dello Unix di Berkeley con una apposita sigla.

Nel 1983, con il rilascio della versione 4.2 di BSD, venne definita una
implementazione delle funzioni di interfaccia a cui si fa riferimento con la
sigla 4.2BSD. Per fare riferimento alle precedenti versioni si usano poi le
sigle 3BSD e 4BSD (per le due versioni pubblicate nel 1980), e 4.1BSD per
quella pubblicata nel 1981.

Le varie estensioni ideate a Berkeley sono state via via aggiunte al sistema
nelle varie versioni succedutesi negli anni, che vanno sotto il nome di
4.3BSD, per la versione rilasciata nel 1986 e 4.4BSD, per la versione
rilasciata nel 1993, che costituisce l'ultima release ufficiale
dell'università di Berkeley. Si tenga presente che molte di queste interfacce
sono presenti in derivati commerciali di BSD come SunOS. Il kernel Linux e la
\acr{glibc} forniscono tutte queste estensioni che sono state in gran parte
incorporate negli standard successivi.


\subsection{Gli standard IEEE -- POSIX}
\label{sec:intro_posix}

Lo standard ufficiale creato da un organismo indipendente più attinente alle
interfacce di un sistema unix-like nel suo complesso (e che concerne sia il
kernel che le librerie che i comandi) è stato lo standard POSIX. Esso prende
origine dallo standard ANSI C, che contiene come sottoinsieme, prevedendo
ulteriori capacità per le funzioni in esso definite, ed aggiungendone di
nuove.

In realtà POSIX è una famiglia di standard diversi, il cui nome, suggerito da
Richard Stallman, sta per \textit{Portable Operating System Interface}, ma la
X finale denuncia la sua stretta relazione con i sistemi Unix. Esso nasce dal
lavoro dell'IEEE (\textit{Institute of Electrical and Electronics Engeneers})
che ne produsse una prima versione, nota come \textsl{IEEE 1003.1-1988},
mirante a standardizzare l'interfaccia con il sistema operativo.

Ma gli standard POSIX non si limitano alla standardizzazione delle funzioni di
libreria, e in seguito sono stati prodotti anche altri standard per la shell e
i comandi di sistema (1003.2), per le estensioni \textit{real-time} e per i
\textit{thread} (rispettivamente 1003.1d e 1003.1c) per i socket (1003.1g) e
vari altri.  In tab.~\ref{tab:intro_posix_std} è riportata una classificazione
sommaria dei principali documenti prodotti, e di come sono identificati fra
IEEE ed ISO; si tenga conto inoltre che molto spesso si usa l'estensione IEEE
anche come aggiunta al nome POSIX; ad esempio è più comune parlare di POSIX.4
come di POSIX.1b.

Si tenga presente inoltre che nuove specifiche e proposte di standardizzazione
si aggiungono continuamente, mentre le versioni precedenti vengono riviste;
talvolta poi i riferimenti cambiano nome, per cui anche solo seguire le
denominazioni usate diventa particolarmente faticoso.

\begin{table}[htb]
  \footnotesize
  \centering
  \begin{tabular}[c]{|l|l|l|l|}
    \hline
    \textbf{Standard} & \textbf{IEEE} & \textbf{ISO} & \textbf{Contenuto} \\
    \hline
    \hline
    POSIX.1 & 1003.1 & 9945-1& Interfacce di base.                          \\
    POSIX.1a& 1003.1a& 9945-1& Estensioni a POSIX.1.                        \\
    POSIX.2 & 1003.2 & 9945-2& Comandi.                                     \\
    POSIX.3 & 2003   &TR13210& Metodi di test.                              \\
    POSIX.4 & 1003.1b &  --- & Estensioni real-time.                        \\
    POSIX.4a& 1003.1c &  --- & Thread.                                      \\
    POSIX.4b& 1003.1d &9945-1& Ulteriori estensioni real-time.              \\
    POSIX.5 & 1003.5  & 14519& Interfaccia per il linguaggio ADA.           \\
    POSIX.6 & 1003.2c,1e& 9945-2& Sicurezza.                                \\
    POSIX.8 & 1003.1f& 9945-1& Accesso ai file via rete.                    \\
    POSIX.9 & 1003.9  &  --- & Interfaccia per il Fortran-77.               \\
    POSIX.12& 1003.1g& 9945-1& Socket.                                      \\
    \hline
  \end{tabular}
  \caption{Elenco dei vari standard POSIX e relative denominazioni.}
  \label{tab:intro_posix_std}
\end{table}

Benché l'insieme degli standard POSIX siano basati sui sistemi Unix, essi
definiscono comunque un'interfaccia di programmazione generica e non fanno
riferimento ad una implementazione specifica (ad esempio esiste
un'implementazione di POSIX.1 anche sotto Windows NT).  

Linux e la \acr{glibc} implementano tutte le funzioni definite nello standard
POSIX.1, queste ultime forniscono in più alcune ulteriori capacità (per
funzioni di \textit{pattern matching} e per la manipolazione delle
\textit{regular expression}), che vengono usate dalla shell e dai comandi di
sistema e che sono definite nello standard POSIX.2.

Nelle versioni più recenti del kernel e delle librerie sono inoltre supportate
ulteriori funzionalità aggiunte dallo standard POSIX.1c per quanto riguarda i
\textit{thread} (vedi cap.~\ref{cha:threads}), e dallo standard POSIX.1b per
quanto riguarda i segnali e lo scheduling real-time
(sez.~\ref{sec:sig_real_time} e sez.~\ref{sec:proc_real_time}), la misura del
tempo, i meccanismi di intercomunicazione (sez.~\ref{sec:ipc_posix}) e l'I/O
asincrono (sez.~\ref{sec:file_asyncronous_io}).

Lo standard principale resta comunque POSIX.1, che continua ad evolversi; la
versione più nota, cui gran parte delle implementazioni fanno riferimento, e
che costituisce una base per molti altri tentativi di standardizzazione, è
stata rilasciata anche come standard internazionale con la sigla
\textsl{ISO/IEC 9945-1:1996} ed include i precedenti POSIX.1b e POSIX.1c. In
genere si fa riferimento ad essa come POSIX.1-1996.

Nel 2001 è stata poi eseguita una sintesi degli standard POSIX.1, POSIX.2 e
SUSv3 (vedi sez.~\ref{sec:intro_xopen}) in un unico documento, redatto sotto
gli auspici del cosiddetto gruppo Austin che va sotto il nome di POSIX.1-2001.
Questo standard definisce due livelli di conformità, quello POSIX, in cui sono
presenti solo le interfacce di base, e quello XSI che richiede la presenza di
una serie di estensioni opzionali per lo standard POSIX, riprese da SUSv3.
Inoltre lo standard è stato allineato allo standard C99, e segue lo stesso
nella definizione delle interfacce.

A questo standard sono stati aggiunti due documenti di correzione e
perfezionamento denominati \textit{Technical Corrigenda}, il TC1 del 2003 ed
il TC2 del 2004, e talvolta si fa riferimento agli stessi con le sigle
POSIX.1-2003 e POSIX.1-2004. 

Una ulteriore revisione degli standard POSIX e SUS è stata completata e
ratificata nel 2008, cosa che ha portato al rilascio di una nuova versione
sotto il nome di POSIX.1-2008 (e SUSv4), con l'incorporazione di alcune nuove
interfacce, la obsolescenza di altre, la trasformazione da opzionali a
richieste di alcune specifiche di base, oltre alle solite precisazioni ed
aggiornamenti. Anche in questo caso è prevista la suddivisione in una
conformità di base, e delle interfacce aggiuntive. Una ulteriore revisione è
stata pubblicata nel 2017 come POSIX.1-2017.

Le procedure di aggiornamento dello standard POSIX prevedono comunque un
percorso continuo, che prevede la possibilità di introduzione di nuove
interfacce e la definizione di precisazioni ed aggiornamenti, per questo in
futuro verranno rilasciate nuove versioni. Alla stesura di queste note
l'ultima revisione approvata resta POSIX.1-2017, uno stato della situazione
corrente del supporto degli standard è allegato alla documentazione della
\acr{glibc} e si può ottenere con il comando \texttt{man standards}.


\subsection{Gli standard X/Open -- Opengroup -- Unix}
\label{sec:intro_xopen}

Il consorzio X/Open nacque nel 1984 come consorzio di venditori di sistemi
Unix per giungere ad un'armonizzazione delle varie implementazioni.  Per far
questo iniziò a pubblicare una serie di documentazioni e specifiche sotto il
nome di \textit{X/Open Portability Guide} a cui di norma si fa riferimento con
l'abbreviazione XPG$n$, con $n$ che indica la versione.

Nel 1989 il consorzio produsse una terza versione di questa guida
particolarmente voluminosa (la \textit{X/Open Portability Guide, Issue 3}),
contenente una dettagliata standardizzazione dell'interfaccia di sistema di
Unix, che venne presa come riferimento da vari produttori. Questo standard,
detto anche XPG3 dal nome della suddetta guida, è sempre basato sullo standard
POSIX.1, ma prevede una serie di funzionalità aggiuntive fra cui le specifiche
delle API (sigla che sta per \textit{Application Programmable Interface}, in
italiano interfaccia di programmazione) per l'interfaccia grafica (X11).

Nel 1992 lo standard venne rivisto con una nuova versione della guida, la
Issue 4, da cui la sigla XPG4, che aggiungeva l'interfaccia XTI (\textit{X
  Transport Interface}) mirante a soppiantare (senza molto successo)
l'interfaccia dei socket derivata da BSD. Una seconda versione della guida fu
rilasciata nel 1994; questa è nota con il nome di Spec 1170 (dal numero delle
interfacce, intestazioni e comandi definiti) ma si fa riferimento ad essa
anche come XPG4v2.

Nel 1993 il marchio Unix passò di proprietà dalla Novell (che a sua volta lo
aveva comprato dalla AT\&T) al consorzio X/Open che iniziò a pubblicare le sue
specifiche sotto il nome di \textit{Single UNIX Specification} o SUS, l'ultima
versione di Spec 1170 diventò così la prima versione delle \textit{Single UNIX
  Specification}, detta SUS o SUSv1, ma più comunemente nota anche come
\textit{Unix 95}.

Nel 1996 la fusione del consorzio X/Open con la Open Software Foundation (nata
da un gruppo di aziende concorrenti rispetto ai fondatori di X/Open) portò
alla costituzione dell'\textit{Open Group}, un consorzio internazionale che
raccoglie produttori, utenti industriali, entità accademiche e governative.
Attualmente il consorzio è detentore del marchio depositato Unix, e prosegue
il lavoro di standardizzazione delle varie implementazioni, rilasciando
periodicamente nuove specifiche e strumenti per la verifica della conformità
alle stesse.

Nel 1997 fu annunciata la seconda versione delle \textit{Single UNIX
  Specification}, nota con la sigla SUSv2, in questa versione le interfacce
specificate salgono a 1434, e addirittura a 3030 se si considerano le stazioni
di lavoro grafiche, per le quali sono inserite pure le interfacce usate da CDE
che richiede sia X11 che Motif. La conformità a questa versione permette l'uso
del nome \textit{Unix 98}, usato spesso anche per riferirsi allo standard. Un
altro nome alternativo di queste specifiche, date le origini, è XPG5.

Come accennato nel 2001, con il rilascio dello standard POSIX.1-2001, è stato
effettuato uno sforzo di sintesi in cui sono state comprese, nella parte di
interfacce estese, anche le interfacce definite nelle \textit{Single UNIX
  Specification}, pertanto si può fare riferimento a detto standard, quando
comprensivo del rispetto delle estensioni XSI, come SUSv3, e fregiarsi del
marchio UNIX 03 se conformi ad esso. 

Infine, come avvenuto per POSIX.1-2001, anche con la successiva revisione
dello standard POSIX.1 (la POSIX.1-2008) è stato stabilito che la conformità
completa a tutte quelle che sono le nuove estensioni XSI previste
dall'aggiornamento vada a definire la quarta versione delle \textit{Single
  UNIX Specification}, chiamata appunto SUSv4.


\subsection{Il controllo di aderenza agli standard}
\label{sec:intro_gcc_glibc_std}

In Linux, se si usa la \acr{glibc}, la conformità agli standard appena
descritti può essere richiesta sia attraverso l'uso di alcune opzioni del
compilatore (il \texttt{gcc}) che con la definizione di opportune costanti
prima dell'inclusione dei file di intestazione (gli \textit{header file}, vedi
sez.~\ref{sec:proc_syscall}) in cui le varie funzioni di libreria vengono
definite.

Ad esempio se si vuole che i programmi seguano una stretta attinenza allo
standard ANSI C si può usare l'opzione \texttt{-ansi} del compilatore, e non
potrà essere utilizzata nessuna funzione non riconosciuta dalle specifiche
standard ISO per il C.  Il \texttt{gcc} possiede inoltre una specifica opzione
per richiedere la conformità ad uno standard, nella forma \texttt{-std=nome},
dove \texttt{nome} può essere \texttt{c89} per indicare lo standard ANSI C
(vedi sez.~\ref{sec:intro_ansiC}) o \texttt{c99} per indicare la conformità
allo standard C99 o \texttt{c11} per indicare la conformità allo standard C11
(revisione del 2011).\footnote{esistono anche le possibilità di usare i valori
  \texttt{gnu89}, che indica l'uso delle estensioni GNU al C89, riprese poi
  dal C99, \texttt{gnu99} che indica il dialetto GNU del C99, o \texttt{gnu11}
  che indica le estensioni GNU al C11, lo standard adottato di default dipende
  dalla versione del \texttt{gcc}, ed all'agosto 2018 con la versione 8.2 è
  \texttt{gnu11}.}

Per attivare le varie opzioni di controllo di aderenza agli standard è poi
possibile definire delle macro di preprocessore che controllano le
funzionalità che la \acr{glibc} può mettere a disposizione:\footnote{le macro
  sono definite nel file di dichiarazione \file{<features.h>}, ma non è
  necessario includerlo nei propri programmi in quanto viene automaticamente
  incluso da tutti gli altri file di dichiarazione che utilizzano le macro in
  esso definite; si tenga conto inoltre che il file definisce anche delle
  ulteriori macro interne, in genere con un doppio prefisso di \texttt{\_},
  che non devono assolutamente mai essere usate direttamente. } questo può
essere fatto attraverso l'opzione \texttt{-D} del compilatore, ma è buona
norma farlo inserendo gli opportuni \code{\#define} prima della inclusione dei
propri \textit{header file} (vedi sez.~\ref{sec:proc_syscall}).

Le macro disponibili per controllare l'aderenza ai vari standard messi a
disposizione della \acr{glibc}, che rendono disponibili soltanto le funzioni
in essi definite, sono illustrate nel seguente elenco:
\begin{basedescript}{\desclabelwidth{2.7cm}\desclabelstyle{\nextlinelabel}}
\item[\macrod{\_\_STRICT\_ANSI\_\_}] richiede l'aderenza stretta allo standard
  C ISO; viene automaticamente predefinita qualora si invochi il \texttt{gcc}
  con le opzione \texttt{-ansi} o \texttt{-std=c99}.

\item[\macrod{\_POSIX\_SOURCE}] definendo questa macro (considerata obsoleta)
  si rendono disponibili tutte le funzionalità dello standard POSIX.1 (la
  versione IEEE Standard 1003.1) insieme a tutte le funzionalità dello
  standard ISO C. Se viene anche definita con un intero positivo la macro
  \macro{\_POSIX\_C\_SOURCE} lo stato di questa non viene preso in
  considerazione.

\item[\macrod{\_POSIX\_C\_SOURCE}] definendo questa macro ad un valore intero
  positivo si controlla quale livello delle funzionalità specificate da POSIX
  viene messa a disposizione; più alto è il valore maggiori sono le
  funzionalità:
  \begin{itemize}
  \item un valore uguale a ``\texttt{1}'' rende disponibili le funzionalità
    specificate nella edizione del 1990 (IEEE Standard 1003.1-1990);
  \item valori maggiori o uguali a ``\texttt{2}'' rendono disponibili le
    funzionalità previste dallo standard POSIX.2 specificate nell'edizione del
    1992 (IEEE Standard 1003.2-1992),
  \item un valore maggiore o uguale a ``\texttt{199309L}'' rende disponibili
    le funzionalità previste dallo standard POSIX.1b specificate nell'edizione
    del 1993 (IEEE Standard 1003.1b-1993);
  \item un valore maggiore o uguale a ``\texttt{199506L}'' rende disponibili
    le funzionalità previste dallo standard POSIX.1 specificate nell'edizione
    del 1996 (\textit{ISO/IEC 9945-1:1996}), ed in particolare le definizioni
    dello standard POSIX.1c per i \textit{thread};
  \item a partire dalla versione 2.3.3 della \acr{glibc} un valore maggiore o
    uguale a ``\texttt{200112L}'' rende disponibili le funzionalità di base
    previste dallo standard POSIX.1-2001, escludendo le estensioni XSI;
  \item a partire dalla versione 2.10 della \acr{glibc} un valore maggiore o
    uguale a ``\texttt{200809L}'' rende disponibili le funzionalità di base
    previste dallo standard POSIX.1-2008, escludendo le estensioni XSI;
  \item in futuro valori superiori potranno abilitare ulteriori estensioni.
  \end{itemize}

\item[\macrod{\_BSD\_SOURCE}] definendo questa macro si rendono disponibili le
  funzionalità derivate da BSD4.3, insieme a quelle previste dagli standard
  ISO C, POSIX.1 e POSIX.2; alcune delle funzionalità previste da BSD sono
  però in conflitto con le corrispondenti definite nello standard POSIX.1, in
  questo caso se la macro è definita le definizioni previste da BSD4.3 avranno
  la precedenza rispetto a POSIX.

  A causa della natura dei conflitti con POSIX per ottenere una piena
  compatibilità con BSD4.3 può essere necessario anche usare una libreria di
  compatibilità, dato che alcune funzioni sono definite in modo diverso. In
  questo caso occorrerà anche usare l'opzione \cmd{-lbsd-compat} con il
  compilatore per indicargli di utilizzare le versioni nella libreria di
  compatibilità prima di quelle normali.

  Si tenga inoltre presente che la preferenza verso le versioni delle funzioni
  usate da BSD viene mantenuta soltanto se nessuna delle ulteriori macro di
  specificazione di standard successivi (vale a dire una fra
  \macro{\_POSIX\_C\_SOURCE}, \macro{\_POSIX\_SOURCE}, \macro{\_SVID\_SOURCE},
  \macro{\_XOPEN\_SOURCE}, \macro{\_XOPEN\_SOURCE\_EXTENDED} o
  \macro{\_GNU\_SOURCE}) è stata a sua volta attivata, nel qual caso queste
  hanno la precedenza. Se però si definisce \macro{\_BSD\_SOURCE} dopo aver
  definito una di queste macro, l'effetto sarà quello di dare la precedenza
  alle funzioni in forma BSD. Questa macro, essendo ricompresa in
  \macro{\_DEFAULT\_SOURCE} che è definita di default, è stata deprecata a
  partire dalla \acr{glibc} 2.20.

\item[\macrod{\_SVID\_SOURCE}] definendo questa macro si rendono disponibili le
  funzionalità derivate da SVID. Esse comprendono anche quelle definite negli
  standard ISO C, POSIX.1, POSIX.2, e X/Open (XPG$n$) illustrati in
  precedenza. Questa macro, essendo ricompresa in \macro{\_DEFAULT\_SOURCE}
  che è definita di default, è stata deprecata a partire dalla \acr{glibc}
  2.20.

\item[\macrod{\_DEFAULT\_SOURCE}] questa macro abilita le definizioni
  considerate il \textit{default}, comprese quelle richieste dallo standard
  POSIX.1-2008, ed è sostanzialmente equivalente all'insieme di
  \macro{\_SVID\_SOURCE}, \macro{\_BSD\_SOURCE} e
  \macro{\_POSIX\_C\_SOURCE}. Essendo predefinita non è necessario usarla a
  meno di non aver richiesto delle definizioni più restrittive sia con altre
  macro che con i flag del compilatore, nel qual caso abilita le funzioni che
  altrimenti sarebbero disabilitate. Questa macro è stata introdotta a partire
  dalla \acr{glibc} 2.19 e consente di deprecare \macro{\_SVID\_SOURCE} e
  \macro{\_BSD\_SOURCE}.

\item[\macrod{\_XOPEN\_SOURCE}] definendo questa macro si rendono disponibili
  le funzionalità descritte nella \textit{X/Open Portability Guide}. Anche
  queste sono un sovrainsieme di quelle definite negli standard POSIX.1 e
  POSIX.2 ed in effetti sia \macro{\_POSIX\_SOURCE} che
  \macro{\_POSIX\_C\_SOURCE} vengono automaticamente definite. Sono incluse
  anche ulteriori funzionalità disponibili in BSD e SVID, più una serie di
  estensioni a secondo dei seguenti valori:
  \begin{itemize}
  \item la definizione della macro ad un valore qualunque attiva le
    funzionalità specificate negli standard POSIX.1, POSIX.2 e XPG4;
  \item un valore di ``\texttt{500}'' o superiore rende disponibili anche le
    funzionalità introdotte con SUSv2, vale a dire la conformità ad Unix98;
  \item a partire dalla versione 2.2 della \acr{glibc} un valore uguale a
    ``\texttt{600}'' o superiore rende disponibili anche le funzionalità
    introdotte con SUSv3, corrispondenti allo standard POSIX.1-2001 più le
    estensioni XSI.
  \item a partire dalla versione 2.10 della \acr{glibc} un valore uguale a
    ``\texttt{700}'' o superiore rende disponibili anche le funzionalità
    introdotte con SUSv4, corrispondenti allo standard POSIX.1-2008 più le
    estensioni XSI.
  \end{itemize}

\item[\macrod{\_XOPEN\_SOURCE\_EXTENDED}] definendo questa macro si rendono
  disponibili le ulteriori funzionalità necessarie la conformità al
  rilascio del marchio \textit{X/Open Unix} corrispondenti allo standard
  Unix95, vale a dire quelle specificate da SUSv1/XPG4v2. Questa macro viene
  definita implicitamente tutte le volte che si imposta
  \macro{\_XOPEN\_SOURCE} ad un valore maggiore o uguale a 500.

\item[\macrod{\_ISOC99\_SOURCE}] definendo questa macro si rendono disponibili
  le funzionalità previste per la revisione delle librerie standard del C
  introdotte con lo standard ISO C99. La macro è definita a partire dalla
  versione 2.1.3 della \acr{glibc}. 

  Le versioni precedenti la serie 2.1.x riconoscevano le stesse estensioni con
  la macro \macro{\_ISOC9X\_SOURCE}, dato che lo standard non era stato
  finalizzato, ma la \acr{glibc} aveva già un'implementazione completa che
  poteva essere attivata definendo questa macro. Benché questa sia obsoleta
  viene tuttora riconosciuta come equivalente di \macro{\_ISOC99\_SOURCE} per
  compatibilità.

\item[\macrod{\_ISOC11\_SOURCE}] definendo questa macro si rendono disponibili
  le funzionalità previste per la revisione delle librerie standard del C
  introdotte con lo standard ISO C11, e abilita anche quelle previste dagli
  standard C99 e C95. La macro è definita a partire dalla versione 2.16 della
  \acr{glibc}.

\item[\macrod{\_GNU\_SOURCE}] definendo questa macro si rendono disponibili
  tutte le funzionalità disponibili nei vari standard oltre a varie estensioni
  specifiche presenti solo nella \acr{glibc} ed in Linux. Gli standard coperti
  sono: ISO C89, ISO C99, POSIX.1, POSIX.2, BSD, SVID, X/Open, SUS.

  L'uso di \macro{\_GNU\_SOURCE} è equivalente alla definizione contemporanea
  delle macro: \macro{\_BSD\_SOURCE}, \macro{\_SVID\_SOURCE},
  \macro{\_POSIX\_SOURCE}, \macro{\_ISOC99\_SOURCE}, e inoltre di
  \macro{\_POSIX\_C\_SOURCE} con valore ``\texttt{200112L}'' (o
  ``\texttt{199506L}'' per le versioni della \acr{glibc} precedenti la 2.5),
  \macro{\_XOPEN\_SOURCE\_EXTENDED} e \macro{\_XOPEN\_SOURCE} con valore 600
  (o 500 per le versioni della \acr{glibc} precedenti la 2.2); oltre a queste
  vengono pure attivate le ulteriori due macro \macro{\_ATFILE\_SOURCE} e
  \macrod{\_LARGEFILE64\_SOURCE} che definiscono funzioni previste
  esclusivamente dalla \acr{glibc}.
 
\end{basedescript}

Benché Linux supporti in maniera estensiva gli standard più diffusi, esistono
comunque delle estensioni e funzionalità specifiche, non presenti in altri
standard e lo stesso vale per la \acr{glibc}, che definisce anche delle
ulteriori funzioni di libreria. Ovviamente l'uso di queste funzionalità deve
essere evitato se si ha a cuore la portabilità, ma qualora questo non sia un
requisito esse possono rivelarsi molto utili.

Come per l'aderenza ai vari standard, le funzionalità aggiuntive possono
essere rese esplicitamente disponibili tramite la definizione di opportune
macro di preprocessore, alcune di queste vengono attivate con la definizione
di \macro{\_GNU\_SOURCE}, mentre altre devono essere attivate esplicitamente,
inoltre alcune estensioni possono essere attivate indipendentemente tramite
una opportuna macro; queste estensioni sono illustrate nel seguente elenco:

\begin{basedescript}{\desclabelwidth{2.7cm}\desclabelstyle{\nextlinelabel}}

\item[\macrod{\_LARGEFILE\_SOURCE}] definendo questa macro si rendono
  disponibili alcune funzioni che consentono di superare una inconsistenza
  presente negli standard con i file di grandi dimensioni, ed in particolare
  definire le due funzioni \func{fseeko} e \func{ftello} che al contrario
  delle corrispettive \func{fseek} e \func{ftell} usano il tipo di dato
  specifico \type{off\_t} (vedi sez.~\ref{sec:file_io}).

\item[\macrod{\_LARGEFILE64\_SOURCE}] definendo questa macro si rendono
  disponibili le funzioni di una interfaccia alternativa al supporto di valori
  a 64 bit nelle funzioni di gestione dei file (non supportati in certi
  sistemi), caratterizzate dal suffisso \texttt{64} aggiunto ai vari nomi di
  tipi di dato e funzioni (come \typed{off64\_t} al posto di \type{off\_t} o
  \funcm{lseek64} al posto di \func{lseek}).

  Le funzioni di questa interfaccia alternativa sono state proposte come una
  estensione ad uso di transizione per le \textit{Single UNIX Specification},
  per consentire la gestione di file di grandi dimensioni anche nei sistemi a
  32 bit, in cui la dimensione massima, espressa con un intero, non poteva
  superare i 2Gb.  Nei nuovi programmi queste funzioni devono essere evitate,
  a favore dell'uso macro \macro{\_FILE\_OFFSET\_BITS}, che definita al valore
  di \texttt{64} consente di usare in maniera trasparente le funzioni
  dell'interfaccia classica.

\item[\macrod{\_FILE\_OFFSET\_BITS}] la definizione di questa macro al valore
  di \texttt{64} consente di attivare la conversione automatica di tutti i
  riferimenti a dati e funzioni a 32 bit nelle funzioni di interfaccia ai file
  con le equivalenti a 64 bit, senza dover utilizzare esplicitamente
  l'interfaccia alternativa appena illustrata. In questo modo diventa
  possibile usare le ordinarie funzioni per effettuare operazioni a 64 bit sui
  file anche su sistemi a 32 bit.\footnote{basterà ricompilare il programma
    dopo averla definita, e saranno usate in modo trasparente le funzioni a 64
    bit.}

  Se la macro non è definita o è definita con valore \texttt{32} questo
  comportamento viene disabilitato, e sui sistemi a 32 bit verranno usate le
  ordinarie funzioni a 32 bit, non avendo più il supporto per file di grandi
  dimensioni. Su sistemi a 64 bit invece, dove il problema non sussiste, la
  macro non ha nessun effetto.

\item[\macrod{\_ATFILE\_SOURCE}] definendo questa macro si rendono disponibili
  le estensioni delle funzioni di creazione, accesso e modifica di file e
  directory che risolvono i problemi di sicurezza insiti nell'uso di
  \textit{pathname} relativi con programmi \textit{multi-thread} illustrate in
  sez.~\ref{sec:file_openat}. Dalla \acr{glibc} 2.10 questa macro viene
  definita implicitamente se si definisce \macro{\_POSIX\_C\_SOURCE} ad un
  valore maggiore o uguale di ``\texttt{200809L}''.

\item[\macrod{\_REENTRANT}] definendo questa macro, o la equivalente
  \macrod{\_THREAD\_SAFE} (fornita per compatibilità) si rendono disponibili
  le versioni rientranti (vedi sez.~\ref{sec:proc_reentrant}) di alcune
  funzioni, necessarie quando si usano i \textit{thread}.  Alcune di queste
  funzioni sono anche previste nello standard POSIX.1c, ma ve ne sono altre
  che sono disponibili soltanto su alcuni sistemi, o specifiche della
  \acr{glibc}, e possono essere utilizzate una volta definita la macro. Oggi
  la macro è obsoleta, già dalla \acr{glibc} 2.3 le funzioni erano già
  completamente rientranti e dalla \acr{glibc} 2.25 la macro è equivalente ad
  definire \macro{\_POSIX\_C\_SOURCE} con un valore di ``\texttt{199606L}''
  mentre se una qualunque delle altre macro che richiede un valore di
  conformità più alto è definita, la sua definizione non ha alcun effetto.

\item[\macrod{\_FORTIFY\_SOURCE}] definendo questa macro viene abilitata
  l'inserimento di alcuni controlli per alcune funzioni di allocazione e
  manipolazione di memoria e stringhe che consentono di rilevare
  automaticamente alcuni errori di \textit{buffer overflow} nell'uso delle
  stesse. La funzionalità è stata introdotta a partire dalla versione 2.3.4
  della \acr{glibc} e richiede anche il supporto da parte del compilatore, che
  è disponibile solo a partire dalla versione 4.0 del \texttt{gcc}.

  Le funzioni di libreria che vengono messe sotto controllo quando questa
  funzionalità viene attivata sono, al momento della stesura di queste note,
  le seguenti: \funcm{memcpy}, \funcm{mempcpy}, \funcm{memmove},
  \funcm{memset}, \funcm{stpcpy}, \funcm{strcpy}, \funcm{strncpy},
  \funcm{strcat}, \funcm{strncat}, \func{sprintf}, \func{snprintf},
  \func{vsprintf}, \func{vsnprintf}, e \func{gets}.

  La macro prevede due valori, con \texttt{1} vengono eseguiti dei controlli
  di base che non cambiano il comportamento dei programmi se si richiede una
  ottimizzazione di livello uno o superiore,\footnote{vale a dire se si usa
    l'opzione \texttt{-O1} o superiore del \texttt{gcc}.}  mentre con il
  valore \texttt{2} vengono aggiunti maggiori controlli. Dato che alcuni dei
  controlli vengono effettuati in fase di compilazione l'uso di questa macro
  richiede anche la collaborazione del compilatore, disponibile dalla
  versione 4.0 del \texttt{gcc}.

\end{basedescript}

Se non è stata specificata esplicitamente nessuna di queste macro il default
assunto è che siano definite \macro{\_BSD\_SOURCE}, \macro{\_SVID\_SOURCE},
\macro{\_POSIX\_SOURCE} e, con le versioni della \acr{glibc} più recenti, che
la macro \macro{\_POSIX\_C\_SOURCE} abbia il valore ``\texttt{200809L}'', per
versioni precedenti della \acr{glibc} il valore assegnato a
\macro{\_POSIX\_C\_SOURCE} era di ``\texttt{200112L}'' prima delle 2.10, di
``\texttt{199506L}'' prima delle 2.4, di ``\texttt{199506L}'' prima delle
2.1. Si ricordi infine che perché queste macro abbiano effetto devono essere
sempre definite prima dell'inclusione dei file di dichiarazione.


% vedi anche man feature_test_macros

% TODO: forse può essere interessante trattare i tunables delle glibc, vedi:
% https://siddhesh.in/posts/the-story-of-tunables.html e
% https://sourceware.org/git/gitweb.cgi?p=glibc.git;a=blob;f=README.tunables;h=0e9b0d7a470aabc2344ad997fca03c7663de0419;hb=HEAD 


% LocalWords:  like kernel multitasking scheduler preemptive sez swap is cap VM
% LocalWords:  everything bootstrap init shell Windows Foundation system call
% LocalWords:  fig libc uClib glibc embedded Library POSIX username PAM Methods
% LocalWords:  Pluggable Autentication group supplementary Name Service Switch
% LocalWords:  LDAP identifier uid gid superuser root if BSD SVr dall' American
% LocalWords:  National Institute International Organisation IEC header tab gcc
% LocalWords:  assert ctype dirent errno fcntl limits malloc setjmp signal utmp
% LocalWords:  stdarg stdio stdlib string times unistd library int short caddr
% LocalWords:  address clock dev ino inode key IPC loff nlink off pid rlim size
% LocalWords:  sigset ssize ptrdiff sys IEEE Richard Portable of TR filesystem
% LocalWords:  Operating Interface dell'IEEE Electrical and Electronics thread
% LocalWords:  Engeneers Socket NT matching regular expression scheduling l'I
% LocalWords:  XPG Portability Issue Application Programmable XTI Transport AT
% LocalWords:  socket Spec Novell Specification SUSv CDE Motif Berkley select
% LocalWords:  SunOS l'AT Sun HP Xenix Description SVID Laboratories MP hsearch
% LocalWords:  drand fmtmsg define SOURCE lbsd compat XOPEN version ISOC Large
% LocalWords:  LARGEFILE Support LFS dell' black rectangle node fill cpu draw
% LocalWords:  ellipse mem anchor west proc SysV SV Definition SCO Austin XSI
% LocalWords:  Technical TC SUS Opengroup features STRICT std ATFILE fseeko VFS
% LocalWords:  ftello fseek ftell lseek FORTIFY REENTRANT SAFE overflow memcpy
% LocalWords:  mempcpy memmove memset stpcpy strcpy strncpy strcat strncat gets
% LocalWords:  sprintf snprintf vsprintf vsnprintf syscall number calendar BITS
% LocalWords:  pathname Google Android standards device Virtual bootloader path
% LocalWords:  filename fifo name components resolution chroot parent symbolic
% LocalWords:  char block VMS raw access MacOS LF CR dos HFS Mac attributes
% LocalWords:  Executable Linkable Format Tool magic descriptor stream locking
% LocalWords:  process mount point ELF

%%% Local Variables: 
%%% mode: latex
%%% TeX-master: "gapil"
%%% End: 
